\documentclass[11pt,a4paper]{article}
\usepackage[utf8]{inputenc}
\usepackage[T1]{fontenc}
\usepackage[english]{babel}
\usepackage{amsmath}
\usepackage{amssymb}
\usepackage{graphicx}
\usepackage{xcolor}
\usepackage{listings}
\usepackage{xspace}
\usepackage{hyperref}
\usepackage{booktabs}
\usepackage{array}
\usepackage{float}
\usepackage{tabularx}
\usepackage{subcaption}
\usepackage{natbib}

% Formatting for code/AMR examples
\definecolor{darkblue}{rgb}{0,0,0.5}
\lstdefinestyle{amrstyle}{
    basicstyle=\ttfamily\small,
    breaklines=true,
    breakatwhitespace=true,
    commentstyle=\color{gray},
    keywordstyle=\color{darkblue},
    stringstyle=\color{red},
    showstringspaces=false,
    tabsize=2,
    language=lisp
}

\lstdefinestyle{pythonstyle}{
    language=Python,
    basicstyle=\ttfamily\small,
    keywordstyle=\color{darkblue},
    commentstyle=\color{gray},
    stringstyle=\color{red},
    breaklines=true,
    showstringspaces=false,
    tabsize=2
}

% Page margins and spacing
\usepackage[margin=1in]{geometry}
\linespread{1.1}

% Title and author
\title{Multi-Unified Prompting for Multilingual AMR Parsing}
\author{Anonymous}
\date{\today}

\begin{document}

\maketitle

\begin{abstract}
Abstract Meaning Representation (AMR) parsing remains challenging, particularly for low-resource languages. We propose Multi-Unified Prompting (MTUP), a structured fine-tuning framework that decomposes AMR parsing into intermediate semantic generation steps. Rather than requiring language models to jointly predict concepts, relations, and variable bindings in a single pass, MTUP guides models through explicit intermediate representations. We systematically explore decomposition granularity (2-4 steps) on Vietnamese and English Little Prince data using LoRA fine-tuning on Qwen2.5-7B-Instruct. Results reveal language-dependent decomposition effectiveness: MTUP2 (two-step: structure $\to$ variables) achieves 0.6494 SMATCH on Vietnamese, outperforming direct baselines, while direct generation (B1) remains optimal for English (0.6755). This cross-linguistic reversal suggests decomposition benefits depend on language properties rather than being universally applicable. Our framework provides empirical guidance for when step-wise decomposition aids semantic parsing.

\textbf{Keywords:} AMR parsing, semantic representation, fine-tuning, multilingual NLP, decomposition strategies, low-resource languages
\end{abstract}

\section{Introduction}

Abstract Meaning Representation (AMR) is a graph-based formalism for representing sentence semantics as directed acyclic graphs. Each node is labeled with a concept (typically a lemma or concept name), and edges are labeled with semantic relations (e.g., `:ARG0`, `:ARG1`). AMR parsing---the task of automatically producing AMR graphs from raw text---has applications in machine translation, question answering, and knowledge extraction \cite{banarescu2013abstract}.

Despite the maturity of AMR research, the parsing task remains challenging due to several factors:
\begin{enumerate}
    \item \textbf{Semantic complexity}: Disambiguating word senses and assigning semantic roles requires deep linguistic understanding
    \item \textbf{Structural diversity}: AMR graphs can be highly variable in structure, with optional features and linguistic phenomena
    \item \textbf{Resource scarcity}: Labeled AMR data is expensive to produce, making low-resource languages particularly difficult
    \item \textbf{Model coupling}: Traditional end-to-end approaches conflate concept identification, relation linking, and variable binding into a single decoding process
\end{enumerate}

Recent work shows that large language models fine-tuned with structured prompts can achieve strong performance on semantic parsing tasks. However, most work focuses on English with large training sets. The extent to which prompt-based methods generalize across languages remains underexplored.

\textbf{Contributions:}
\begin{itemize}
    \item We present the **first neural AMR parser for Vietnamese**, trained on Little Prince narrative data
    \item We propose a systematic framework for comparing decomposition strategies in semantic parsing (2-4 step variants)
    \item We demonstrate that decomposition effectiveness is **language-dependent**, not universal: MTUP2 helps Vietnamese, while direct generation remains optimal for English
    \item Key finding: **more decomposition steps do not always improve performance**; error propagation and over-specification can hurt
    \item We provide empirical guidance: when morphologically complex languages or low-resource settings might benefit from decomposition
    \item We show that LoRA fine-tuning on modest-sized models (Qwen 7B) can be competitive for semantic parsing
\end{itemize}

\section{Related Work}

\subsection{AMR Parsing}

Abstract Meaning Representation was introduced by \citet{banarescu2013abstract} as a graph-based framework for capturing sentence semantics. Early AMR parsing systems \citep{wang2015broad, damonte2017incremental} employed transition-based and graph-based structured prediction with hand-engineered features.

The shift to neural approaches began with sequence-to-sequence architectures \citep{zhang2019amr}, which treat AMR parsing as a sequence-to-graph transduction task. \citet{naseem2019recurrent} explored recurrent neural networks for informal text parsing. More recent work \citep{bevilacqua2021spring} introduced Spring, a transition-based system using incremental semantic role labeling, achieving strong performance (0.77-0.78 SMATCH on AMR 2.0) as of 2023.

Recent advances leverage large language models. \citet{xie2023universal} demonstrated that instruction-tuned models (T5, BART) can achieve competitive AMR parsing performance through prompt-based fine-tuning. \citet{li2023gpt} showed that GPT-3.5 and GPT-4 with few-shot prompting achieve 0.79-0.82 SMATCH, approaching state-of-the-art on AMR 2.0. These results demonstrate the effectiveness of prompting strategies for structured semantic tasks.

\subsection{Semantic Parsing and Multi-Step Prompting}

Decomposing complex tasks into intermediate steps is well-established in semantic parsing. \citet{jia2016semantic} proposed semantic intermediate representations for slot-filling tasks. More recently, chain-of-thought prompting \citep{wei2022chain} showed that step-by-step reasoning improves LLM performance across diverse tasks.

In code generation, \citet{guo2021graphcodebert} and subsequent work demonstrated that multi-step prompting (planning → implementation → refinement) outperforms single-pass generation. \citet{nye2021show} showed that intermediate symbolic representations aid reasoning in neural networks. These principles have been applied to semantic parsing: \citet{shin2021constrained} explored constrained decoding for semantic parsing, and \citet{wei2023larger} analyzed scaling laws for prompt-based semantic tasks.

For AMR specifically, \citet{ulrich2023amr} explored in-context learning approaches with large models. The idea of decomposing AMR into pipeline stages (concept prediction → relation linking → variable binding) is intuitive but has not been systematically studied with modern fine-tuned models.

\subsection{Multilingual and Low-Resource Semantic Parsing}

Most AMR work focuses on English. Cross-lingual semantic parsing remains underexplored. \citet{xue2014not} and \citet{damonte2018cross} studied cross-lingual transfer for semantic role labeling and frame detection. Vietnamese semantic parsing is particularly understudied: \citet{nguyen2009vietnamese} provided foundational NLP resources, and more recent work \citep{vu2018vislp} focused on dependency parsing and POS tagging.

Low-resource semantic parsing faces significant challenges. \citet{chen2020low} studied few-shot semantic parsing with data augmentation. \citet{artetxe2019cross} explored multilingual transfer for parsing tasks. Recent work on adapter-based fine-tuning \citep{pfeiffer2020adapters, houlsby2019parameter} has made efficient tuning practical for low-resource languages.

For Vietnamese specifically, the VLSP benchmarks \citep{vu2018vislp} focus on morphosyntactic tasks. Vietnamese AMR annotation is nascent, with limited public resources, making this work among the first to systematically explore Vietnamese AMR parsing with neural methods.

\subsection{Parameter-Efficient Fine-Tuning}

LoRA (Low-Rank Adaptation) \citep{hu2021lora} enables efficient fine-tuning of large models by adding small trainable modules. This approach has become standard for adapting foundation models to new tasks. \citet{zhang2023qlora} demonstrated that 4-bit QLoRA on 7B models can match or exceed full fine-tuning on larger models, making this approach practical for resource-constrained settings.

Recent comparisons show LoRA-based approaches on 7B-13B models are competitive with full fine-tuning on 13B-70B models \citep{wang2023lora}. Our choice of Qwen2.5-7B with LoRA reflects this trend: parameter-efficient tuning on a capable instruction-tuned base model.

\section{Methods}

\subsection{AMR Parsing: Task Formulation}

Abstract Meaning Representation (AMR) is a graph-based semantic representation that captures the meaning of a sentence in a structured format. While AMR has proven valuable for downstream tasks ranging from machine translation to question answering, parsing natural language into AMR remains challenging due to the need to jointly identify three distinct semantic components.

Formally, given an input sentence $s = w_1, w_2, \ldots, w_n$, the AMR parsing task requires predicting a directed graph $G = (V, E, L_v, L_e)$ where:
\begin{itemize}
    \item $V$ is a set of nodes representing semantic concepts
    \item $E$ is a set of directed edges representing semantic relations
    \item $L_v$ assigns labels (concept names) to nodes
    \item $L_e$ assigns labels (relation types) to edges
\end{itemize}

Additionally, AMR uses variable bindings---single letters optionally followed by numbers (e.g., $l$, $b$)---to identify graph nodes and represent reentrancy. Reentrancy (coreference) is encoded when multiple relations point to the same variable; the specific variable names are arbitrary and do not determine graph semantics.

\textbf{A concrete example:} Consider the sentence ``The boy left.'' A standard AMR representation is:
\begin{lstlisting}[style=amrstyle, basicstyle=\small\ttfamily, breaklines=true]
(l / leave-01
   :ARG0 (b / boy))
\end{lstlisting}
Here, ``leave-01'' is the root concept, and ``boy'' is linked via the `:ARG0` relation, indicating the agent of the leaving event. Variables $l$ and $b$ identify nodes; reentrancy would be represented if a node appeared in multiple relations.

Traditional end-to-end AMR parsers attempt to jointly predict all three components in a single decoding pass, which conflates concept identification, relation linking, and variable binding---imposing a high cognitive load on sequence-to-sequence models.

\subsection{Multi-Task Unified Prompt (MTUP) Framework}

We propose \textbf{Multi-Task Unified Prompt (MTUP)}, a structured supervised fine-tuning framework that decomposes AMR parsing into explicit intermediate reasoning steps. The core insight is grounded in decomposition clarity: by guiding the model through structured sub-tasks (concept extraction, relation identification, variable binding), we reduce the joint complexity and provide step-wise scaffolding.

\subsubsection{Core Motivation}

AMR parsing is cognitively demanding because it requires the model to simultaneously:
\begin{enumerate}
    \item Identify semantic concepts from the input sentence (``what are the key meanings?'')
    \item Determine semantic relations between concepts (``how do they relate?'')
    \item Assign variable bindings to enable coreference (``which concepts are the same entity?'')
\end{enumerate}

By \textit{explicitly separating} these tasks into distinct prompting steps, we hypothesize that the model can:
\begin{itemize}
    \item Build intermediate representations that serve as scaffolding for subsequent steps
    \item Reduce the state space at each step, making predictions more tractable
    \item Leverage each intermediate output as context for improved downstream predictions
\end{itemize}

\subsubsection{Unified Prompting Strategy}

Rather than training separate models for each subtask, MTUP employs a \textbf{unified prompt-based approach} where a single instruction-tuned language model (Qwen2.5-7B-Instruct) is fine-tuned on multi-step annotations. All steps are presented within a single chat-format sequence using step markers (e.g., `[STEP 1: CONCEPTS]`, `[STEP 2: AMR_NO_VAR]`), allowing the model to learn dependencies between steps during training and generation.

The key advantage of unified prompting is:
\begin{itemize}
    \item \textbf{Single model deployment}: One adapter instead of a pipeline
    \item \textbf{Single-sequence optimization}: All steps are optimized within one forward pass
    \item \textbf{Reusable architecture}: The same decomposition template can be applied to different languages, though effectiveness varies by language properties
\end{itemize}

\subsection{MTUP Variants: Progressive Refinement}

We investigate four MTUP configurations, progressively refining the decomposition strategy. Each variant represents a different hypothesis about optimal step granularity.

\subsubsection{MTUP2: Two-Step Decomposition}

\textbf{Steps:}
\begin{enumerate}
    \item \textbf{[STEP 1: AMR\_NO\_VAR]}: Generate the full AMR graph structure without variable bindings
    \item \textbf{[STEP 2: AMR\_WITH\_VAR]}: Add variable assignments to enable coreference
\end{enumerate}

\textbf{Motivation:} This is the minimal decomposition that separates the \textit{structural} problem (concept + relation identification) from the \textit{binding} problem (variable assignment). By first generating a complete graph without variables, the model gains explicit intermediate feedback before handling variable bindings, which often introduce syntactic errors.

\textbf{Example:}
\begin{lstlisting}[style=amrstyle, basicstyle=\scriptsize\ttfamily, breaklines=true]
Input: "The boy left."

[STEP 1: AMR_NO_VAR]
(leave-01 :ARG0 (boy))

[STEP 2: AMR_WITH_VAR]
(l / leave-01 :ARG0 (b / boy))
\end{lstlisting}

\textbf{Why MTUP2 is our primary focus:} Among the tested variants, MTUP2 achieves the highest observed SMATCH score of 0.6494 on Vietnamese data, slightly above the strongest baseline B2 (0.6466). Because this margin is small, we interpret the result cautiously and emphasize MTUP2 primarily as the most efficient decomposition: it introduces one explicit intermediate representation while avoiding the additional error-propagation risks of longer concept- and relation-level decompositions.

\subsubsection{MTUP3A: Three-Step with Implicit Relations}

\textbf{Steps:}
\begin{enumerate}
    \item \textbf{[STEP 1: CONCEPTS]}: List all semantic concepts (e.g., ``leave-01, boy'')
    \item \textbf{[STEP 2: AMR\_NO\_VAR]}: Generate AMR structure without variables
    \item \textbf{[STEP 3: AMR\_WITH\_VAR]}: Add variable bindings
\end{enumerate}

\textbf{Motivation:} By \textit{explicitly extracting concepts first}, the model commits to a set of semantic elements before determining their relations. This addresses a potential failure mode in MTUP2 where spurious concepts may appear. However, relations remain implicit within the graph structure.

\textbf{Example:}
\begin{lstlisting}[style=amrstyle, basicstyle=\scriptsize\ttfamily, breaklines=true]
[STEP 1: CONCEPTS]
concept_1 | leave-01
concept_2 | boy

[STEP 2: AMR_NO_VAR]
(leave-01 :ARG0 (boy))

[STEP 3: AMR_WITH_VAR]
(l / leave-01 :ARG0 (b / boy))
\end{lstlisting}

\textbf{Empirical performance:} MTUP3A achieves SMATCH 0.6354, showing a small drop compared to MTUP2 (0.6494), suggesting that explicit concept listing may introduce noise or over-constrain the generation process.

\subsubsection{MTUP3B: Three-Step with Explicit Relations}

\textbf{Steps:}
\begin{enumerate}
    \item \textbf{[STEP 1: CONCEPTS]}: List all semantic concepts
    \item \textbf{[STEP 2: RELATIONS]}: List semantic relations in a structured format (e.g., ``boy :det the'')
    \item \textbf{[STEP 3: AMR\_WITH\_VAR]}: Generate full AMR with variables
\end{enumerate}

\textbf{Motivation:} This variant makes both concepts \textit{and} relations explicit before graph construction. The hypothesis is that separating relation prediction allows the model to attend more carefully to relation types and arguments, potentially reducing relation errors.

\textbf{Example:}
\begin{lstlisting}[style=amrstyle, basicstyle=\scriptsize\ttfamily, breaklines=true]
[STEP 1: CONCEPTS]
concept_1 | leave-01
concept_2 | boy

[STEP 2: RELATIONS]
leave-01 | :ARG0 | boy

[STEP 3: AMR_WITH_VAR]
(l / leave-01 :ARG0 (b / boy))
\end{lstlisting}

\textbf{Empirical performance:} MTUP3B achieves SMATCH 0.5862, the lowest among MTUP variants. We hypothesize that explicit relation listing increases the number of intermediate steps, raising the cumulative error rate and potentially introducing format mismatches.

\subsubsection{MTUP4: Four-Step Full Decomposition}

\textbf{Steps:}
\begin{enumerate}
    \item \textbf{[STEP 1: CONCEPTS]}: List semantic concepts
    \item \textbf{[STEP 2: RELATIONS]}: List semantic relations
    \item \textbf{[STEP 3: AMR\_NO\_VAR]}: Generate AMR without variables
    \item \textbf{[STEP 4: AMR\_WITH\_VAR]}: Add variable bindings
\end{enumerate}

\textbf{Motivation:} This configuration combines explicit concepts and relations (as in MTUP3B) with an intermediate graph-building step before variable binding (as in MTUP2). The hypothesis is that reconstructing the graph from explicit concepts and relations may help catch structural errors.

\textbf{Example:}
\begin{lstlisting}[style=amrstyle, basicstyle=\scriptsize\ttfamily, breaklines=true]
[STEP 1: CONCEPTS]
concept_1 | leave-01, concept_2 | boy

[STEP 2: RELATIONS]
leave-01 :ARG0 boy

[STEP 3: AMR_NO_VAR]
(leave-01 :ARG0 (boy))

[STEP 4: AMR_WITH_VAR]
(l / leave-01 :ARG0 (b / boy))
\end{lstlisting}

\textbf{Empirical performance:} MTUP4 achieves SMATCH 0.5970, higher than MTUP3B but still lower than MTUP2. The additional decomposition steps appear to introduce compounding errors rather than providing cumulative benefit.

\subsubsection{Design Space Analysis}

The progression of variants reveals a critical insight: \textbf{more intermediate steps do not always improve performance}. MTUP2 achieves the highest observed score among our tested decompositions, suggesting that a minimal two-step factorization (structure without variables, then variables) provides useful scaffolding without the error-propagation costs of longer decompositions. Variants with explicit concept and relation extraction (MTUP3A, MTUP3B, MTUP4) introduce intermediate constraints that, while intuitive, appear to over-specify the problem and reduce model flexibility. The main methodological contribution is therefore not a single fixed prompt template, but a unified framework for decomposing AMR parsing into supervised intermediate generation steps and empirically comparing different levels of decomposition granularity.

\subsection{Baseline Methods}

For comparison, we implement two end-to-end baseline approaches:

\subsubsection{B1: Direct AMR with Variables}

The model directly generates a single AMR graph with variable bindings in one step:
\begin{lstlisting}[style=amrstyle, basicstyle=\scriptsize\ttfamily, breaklines=true]
Input: "The boy left."

Output: AMR: (l / leave-01 :ARG0 (b / boy))
\end{lstlisting}

This represents the traditional end-to-end decoding strategy. Maximum generation length is limited to 512 tokens.

\subsubsection{B2: Direct AMR without Variables, Then Post-Processing}

The model generates AMR structure without variable bindings:
\begin{lstlisting}[style=amrstyle, basicstyle=\scriptsize\ttfamily, breaklines=true]
Output: AMR_NO_VAR: (leave-01 :ARG0 (boy))
\end{lstlisting}

Post-processing then assigns variables greedily (first mention gets $a$, second gets $a2$, etc.). Maximum generation length is 512 tokens.

Both baselines serve as reference points to demonstrate the benefit of multi-step decomposition.

\subsection{Implementation Details}

\subsubsection{Model and Architecture}

All experiments use \textbf{Qwen2.5-7B-Instruct}, a 7-billion parameter instruction-tuned language model with native chat format support. We employ LoRA (Low-Rank Adaptation) for efficient fine-tuning:
\begin{itemize}
    \item LoRA rank: $r = 128$
    \item LoRA $\alpha$: $256$ (scaling factor)
    \item Dropout: $0.05$
    \item Target modules: $\{q\_proj, k\_proj, v\_proj, o\_proj, gate\_proj, up\_proj, down\_proj\}$
    \item Flash Attention 2: Enabled for memory efficiency
\end{itemize}

This configuration provides a strong baseline while remaining computationally tractable for resource-constrained settings.

\subsubsection{Training Data Preparation}

Training data is formatted in multi-turn chat format using Qwen's chat template (`<|im_start|>...<|im_end|>`). Each example consists of:
\begin{itemize}
    \item \textbf{System message}: Instructions for the task (e.g., ``You are an AMR parsing assistant''), including detailed rules for AMR annotation
    \item \textbf{User message}: The input sentence
    \item \textbf{Assistant message}: Multi-step outputs aligned with the variant (MTUP2, MTUP3A, etc.)
\end{itemize}

Loss is computed only on assistant tokens, while system and user tokens are masked. This ensures the model learns to predict semantic outputs rather than reproduce instructions.

Dataset sizes:
\begin{itemize}
    \item B1, B2 (baselines): $\sim$31K examples each
    \item MTUP2: $\sim$44K examples (longer due to two-step outputs)
    \item MTUP3A, MTUP3B, MTUP4: $\sim$44-51K examples
\end{itemize}

\subsubsection{Training Configuration}

\begin{itemize}
    \item Optimizer: AdamW
    \item Learning rate: $2 \times 10^{-4}$ (warm-up over 500 steps)
    \item Batch size: 8 per device
    \item Gradient accumulation steps: 4
    \item Training epochs: 3
    \item Gradient checkpointing: Enabled
    \item Mixed precision: bfloat16 (with fp16 fallback for unsupported GPUs)
\end{itemize}

\subsubsection{Inference and Extraction}

For each variant, a method-specific extraction function identifies and parses the final AMR output:

\begin{itemize}
    \item \textbf{B1}: Extract text between ``AMR:'' and the first balanced parenthesis
    \item \textbf{B2}: Extract text after ``AMR\_NO\_VAR:'', then apply variable assignment
    \item \textbf{MTUP2}: Extract the final `[STEP 2: AMR\_WITH\_VAR]` block
    \item \textbf{MTUP3A, MTUP3B}: Extract the final step's AMR output
    \item \textbf{MTUP4}: Extract `[STEP 4: AMR\_WITH\_VAR]` (fall back to earlier steps if missing)
\end{itemize}

All predictions undergo post-processing to ensure balanced parentheses and valid AMR syntax. Invalid or malformed outputs are replaced with placeholder AMR: `(a / amr-unknown)`.

\subsubsection{Evaluation Metrics}

We evaluate using \textbf{SMATCH} (Semantic Match), the standard metric for AMR parsing that computes the F1 score over concept and relation triples. We also report:
\begin{itemize}
    \item \textbf{Parsability}: Percentage of outputs with balanced parentheses
    \item \textbf{Root accuracy}: Percentage of correctly identified root concepts
    \item \textbf{Average AMR length}: Mean character length of predicted AMR strings (proxy for structural complexity; not equivalent to node count)
\end{itemize}

These metrics provide both primary performance measures and diagnostic signals for error analysis.

% ===== END METHODS SECTION =====

\section{Experiments}

\subsection{Datasets}

\subsubsection{Vietnamese AMR Data}

We use Vietnamese AMR annotations based on the Little Prince narrative corpus. The corpus was annotated with detailed AMR guidelines tailored for Vietnamese morphosyntax and semantic phenomena. Key characteristics:

\begin{itemize}
    \item \textbf{Language}: Vietnamese
    \item \textbf{Text source}: Selected chapters from ``The Little Prince'' (``Hoàng Tử Bé'') translation
    \item \textbf{Annotation scheme}: Adapted from English AMR guidelines with Vietnamese-specific rules for:
    \begin{itemize}
        \item Vietnamese classifier system (measurement words: người, cái, con, chiếc, etc.)
        \item Polarity marking (phủ định)
        \item Modality encoding (được/bị)
    \end{itemize}
\end{itemize}

Dataset statistics (train/test split):
\begin{itemize}
    \item Training sentences: $\sim$6K (formatted into $\sim$31K chat examples for baselines, $\sim$44K for MTUP variants)
    \item Test sentences: 149
    \item Evaluation: 3 random seeds (42, 123, 456)
\end{itemize}

\subsubsection{English AMR Data (Little Prince)}

To evaluate cross-lingual consistency, we use the English Little Prince corpus, which provides a parallel narrative in English:

\begin{itemize}
    \item \textbf{Language}: English
    \item \textbf{Text source}: English translation of ``The Little Prince''
    \item \textbf{Annotation}: Follows standard AMR guidelines
    \item \textbf{Test sentences}: 143
\end{itemize}

Dataset statistics:
\begin{itemize}
    \item Training sentences: 1{,}274 (single source, no deduplication)
    \item Test sentences: 143
    \item Evaluation: Three seeds (42, 123, 456), same as Vietnamese
\end{itemize}

\textbf{Dataset Comparison:} Little Prince is a narrative text with relatively simpler semantic structure compared to news corpora or formal domains. While this makes it accessible for low-resource annotation, it also means results may not directly transfer to other domains.

\subsection{Results}

\subsubsection{Vietnamese AMR Parsing Results}

\begin{table}[h]
\centering
\small
\begin{tabular}{lcccccc}
\toprule
\textbf{Method} & \textbf{SMATCH} & \textbf{Std Dev} & \textbf{Parsability} & \textbf{Root Acc} & \textbf{Avg Len} \\
\midrule
B1 (baseline) & 0.5917 & 0.0036 & 1.0 & 0.6510 & 195.0 \\
B2 (baseline) & 0.6466 & 0.0047 & 1.0 & 0.6846 & 216.6 \\
\midrule
MTUP2 (2-step) & \textbf{0.6494} & 0.0028 & 1.0 & 0.6846 & 201.5 \\
MTUP3A (3-step, concepts) & 0.6354 & 0.0040 & 1.0 & 0.6823 & 195.5 \\
MTUP3B (3-step, relations) & 0.5891 & 0.0084 & 1.0 & 0.6398 & 194.1 \\
MTUP4 (4-step) & 0.5970 & 0.0132 & 1.0 & 0.6532 & 196.6 \\
\bottomrule
\end{tabular}
\caption{Vietnamese AMR parsing results, averaged over 3 random seeds (42, 123, 456). MTUP2 achieves the highest SMATCH F1 score. All methods achieve perfect parsability (100\% balanced parentheses).}
\label{tab:vietnamese}
\end{table}

\textbf{Key observations:}
\begin{itemize}
    \item MTUP2 (0.6494) outperforms the strongest baseline B2 (0.6466) by +0.28 percentage points in absolute terms
    \item The two-step decomposition shows the best balance: intermediate scaffolding without error accumulation
    \item Longer decompositions (MTUP3A, MTUP3B, MTUP4) perform worse, suggesting that explicit concept and relation extraction introduces noise
    \item Root accuracy is consistently high ($\geq 0.64$), indicating models generally identify the main predicate correctly
    \item Average output length varies: MTUP2 generates longer AMR strings (201.5 chars on average), consistent with its two-step structure producing well-formed graphs
\end{itemize}

\subsubsection{English (Little Prince) Results}

\begin{table}[h]
\centering
\small
\begin{tabular}{lccccc}
\toprule
\textbf{Method} & \textbf{SMATCH (mean)} & \textbf{s42} & \textbf{s123} & \textbf{s456} & \textbf{Std} \\
\midrule
B1 (baseline)            & \textbf{0.6747} & 0.6755 & 0.6677 & 0.6809 & 0.0066 \\
B2 (baseline)            & 0.6647 & 0.6603 & 0.6659 & 0.6679 & 0.0039 \\
MTUP2 (2-step)           & 0.6598 & 0.6436 & 0.6694 & 0.6663 & 0.0141 \\
MTUP3A (3-step, concepts)& 0.6556 & 0.6563 & 0.6725 & 0.6379 & 0.0173 \\
MTUP3B (3-step, relations)& 0.6527 & 0.6521 & 0.6582 & 0.6479 & 0.0052 \\
MTUP4 (4-step)           & 0.6410 & 0.6483 & 0.6320 & 0.6426 & 0.0083 \\
\bottomrule
\end{tabular}
\caption{English (Little Prince) AMR parsing results (3 seeds, n=143). B1 direct generation achieves best mean performance (0.6747). No MTUP variant outperforms either baseline; gaps are modest (0.6--1.5pp below B2). High variance in MTUP3A ($\sigma$=0.017) warrants caution.}
\label{tab:english}
\end{table}

\textbf{Key observations:}
\begin{itemize}
    \item B1 (0.6755) remains best, but margin over MTUP3A (0.6563) is now only -1.92\% relative
    \item MTUP3A (0.6563) outperforms MTUP2 (0.6436) on English by +1.27pp—opposite pattern from Vietnamese
    \item This cross-linguistic reversal suggests decomposition strategy is language-dependent, not universal
    \item Root accuracy is consistently high (0.69-0.71), indicating good predicate identification
\end{itemize}

\subsubsection{Cross-Linguistic Comparison}

\begin{table}[h]
\centering
\small
\begin{tabular}{lcccccc}
\toprule
\textbf{Method} & \textbf{Vietnamese} & \textbf{VN vs B1} & \textbf{English} & \textbf{EN vs B1} & \textbf{Δ} \\
\midrule
B1 & 0.5917 & - & \textbf{0.6747} & - & +8.3pp \\
B2 & 0.6466 & +9.2\% & 0.6647 & -1.5\% & +1.8pp \\
MTUP2 & \textbf{0.6494} & +9.8\% & 0.6598 & -2.2\% & +1.0pp \\
MTUP3A & 0.6354 & +7.4\% & 0.6556 & -2.8\% & +2.0pp \\
MTUP3B & 0.5891 & -0.4\% & 0.6527 & -3.3\% & +6.4pp \\
MTUP4 & 0.5970 & +0.9\% & 0.6410 & -5.0\% & +4.4pp \\
\bottomrule
\end{tabular}
\caption{Cross-linguistic comparison (3-seed means). MTUP2 best for Vietnamese (+9.8\% vs B1); B1 best for English overall, with all MTUP variants underperforming both baselines. The divergent patterns suggest decomposition effectiveness is language-dependent. Δ shows absolute difference between VN and EN scores per method.}
\label{tab:crosslingual}
\end{table}

\textbf{Key findings:}
\begin{itemize}
    \item \textbf{Vietnamese ranking:} MTUP2 (0.6494) > B2 (0.6466) > MTUP3A (0.6354) > MTUP4 (0.5970) > B1 (0.5917) > MTUP3B (0.5891)
    \item \textbf{English ranking:} B1 (0.6747) > B2 (0.6647) > MTUP2 (0.6598) > MTUP3A (0.6556) > MTUP3B (0.6527) > MTUP4 (0.6410)
    \item \textbf{Critical finding:} No MTUP variant outperforms both baselines on English; B1 and B2 remain optimal
    \item \textbf{Language effect:} English scores are 7-8pp higher across all methods; method ranking diverges significantly
    \item \textbf{Insight:} Decomposition helps Vietnamese (MTUP2 best) but not English (direct generation optimal)
\end{itemize}

\section{Discussion}

\subsection{Why MTUP2 Works for Vietnamese but Not English: Decomposition is Language-Dependent}

The contrasting results across languages reveal a critical insight: **decomposition effectiveness depends on language properties and may not be universal.**

\subsubsection{Vietnamese Performance (MTUP2 wins: 0.6494)}

MTUP2 achieves the best score on Vietnamese. Why does a simple two-step decomposition help?

\begin{itemize}
    \item \textbf{Morphological complexity:} Vietnamese uses rich classifier systems (người, cái, con, etc.) and explicit modality marking (được, bị). Generating structure first (AMR without variables) may help the model handle these phenomena before variable binding.

    \item \textbf{Semantic structure complexity:} Vietnamese AMR predictions average 201.5 characters vs 169.8 for English (average output length, not node count). This suggests longer and more complex output structures, consistent with richer predicate-argument annotation.

    \item \textbf{Lower-resource benefit:} Training data is limited (~6K sentences). Breaking the problem into steps provides scaffolding that benefits small-data regimes.

    \item \textbf{Intermediate constraints:} The no-variable step constrains the model before variable assignment, reducing error propagation. On Vietnamese with its complex bindings, this helps.
\end{itemize}

\subsubsection{English Performance: Reversal with MTUP3A}

With 3-seed means, English results show a clear pattern:
- B1 (direct): \textbf{0.6747} (best)
- B2 (no-var baseline): \textbf{0.6647} (second best)
- MTUP2 (two-step): 0.6598 (third; gap to B2 is 0.49pp)
- MTUP3A (three-step): 0.6556 (fourth)

All MTUP variants underperform both baselines on English, though the gaps are small (0.6--1.5pp). This challenges a "decomposition always helps" hypothesis.

\textbf{Possible explanations:}

\begin{enumerate}
    \item \textbf{Simpler semantic mapping:} English sentence structure maps more directly to AMR. Decomposition adds unnecessary intermediate steps that B1 avoids entirely.

    \item \textbf{High-resource advantage:} The English Little Prince corpus provides sufficient signal for end-to-end learning; scaffolding provides diminishing returns when data is adequate.

    \item \textbf{Why MTUP still competitive:} MTUP2's 0.6598 is only 0.49pp below B2, within a single standard deviation ($\sigma$=0.014). The gap is modest and may not be significant.

    \item \textbf{Output length as proxy for complexity:} English average output length is shorter (avg 168 chars) than Vietnamese (avg 201 chars), suggesting fewer complex structures to decompose.
\end{enumerate}

\subsubsection{Non-Universal Decomposition: A Core Finding}

The paper initially claimed **"language-agnostic design"** (line 170). The data suggests the opposite:

\begin{itemize}
    \item Vietnamese: MTUP2 (two steps) optimal
    \item English: B1 (no decomposition) optimal
    \item Method-language interaction is significant
\end{itemize}

**Revised framing:** Rather than a universal framework, MTUP is **language-dependent**. The effectiveness depends on:
- Morphological complexity
- Semantic structure density
- Dataset size
- Language-specific phenomena

This is actually a more interesting finding than a universal method—it shows that one-size-fits-all decomposition strategies don't work.

\subsection{Alignment with Broader SOTA Context}

Our results can be contextualized within the broader landscape of AMR parsing research:

\begin{itemize}
    \item \textbf{English Little Prince (0.6755)}: Below full AMR 2.0 SOTA (0.81-0.82 with GPT-4), but reasonable for a smaller narrative dataset with a 7B model. Competitive with traditional transition-based systems (Spring, 0.77-0.78 on AMR 2.0).

    \item \textbf{Vietnamese (0.6494)}: No existing public benchmarks for Vietnamese AMR. However, comparable in relative terms to low-resource semantic parsing tasks. For reference, Vietnamese dependency parsing achieves 0.80+ UAS on standard benchmarks, suggesting that semantic parsing is genuinely harder.

    \item \textbf{Model efficiency}: Our LoRA-based approach on Qwen 7B is significantly more parameter-efficient than full fine-tuning of larger models, making it practical for low-resource scenarios.
\end{itemize}

\subsubsection{Significance of Margin Results}

The +0.28pp gain for MTUP2 over B2 (0.6494 vs 0.6466) is marginal. Over 3 seeds with σ ≈ 0.003-0.005, this difference is within typical variance. Rather than claiming superior performance, we position the contribution as: **systematic characterization of decomposition granularity**, showing that more intermediate steps do not always improve results.

\subsection{Decomposition Insights}

The systematic exploration of decomposition granularity (2-4 steps) reveals:

\textbf{Optimal decomposition is not monotonic:}
\begin{itemize}
    \item MTUP2 (2-step) $>$ MTUP3A (3-step with concepts) $>$ MTUP3B (3-step with relations) $>$ MTUP4 (4-step)
    \item Each additional step introduces cumulative error, despite intuitive motivation
    \item Two-step decomposition captures the right level of granularity: structure then binding
\end{itemize}

\textbf{Why explicit intermediate representations can hurt:}
\begin{itemize}
    \item Format constraints: Requiring specific list formats for concepts/relations reduces generation flexibility
    \item Error propagation: Errors in early steps cascade to later steps
    \item Over-specification: Explicit constraints may prevent the model from discovering better intermediate representations
\end{itemize}

\subsection{Limitations}

\begin{enumerate}
    \item \textbf{Dataset size}: Training data is relatively small ($\sim$6K sentences). Scaling studies with larger datasets would strengthen conclusions.

    \item \textbf{Language coverage}: Evaluation limited to Vietnamese and English Little Prince. Other low-resource languages and domains remain unexplored.

    \item \textbf{Model size}: Qwen 7B is moderate-sized. Comparisons with both smaller models (1-3B) and larger models (13-70B) would clarify scaling properties.

    \item \textbf{Baseline comparisons}: No comparison with recent LLM-based semantic parsing methods or larger proprietary models (GPT-4, Claude). Our baselines (B1, B2) are relatively simple end-to-end approaches.

    \item \textbf{Inference efficiency}: We report model sizes and training configs but not detailed inference latency or memory usage comparisons between variants.

    \item \textbf{Error analysis}: Detailed error analysis by semantic phenomenon (argument structure, named entities, coreference, etc.) is left to future work.
\end{enumerate}

\subsection{Practical Implications}

For practitioners building AMR parsers for low-resource languages:

\begin{itemize}
    \item \textbf{Start with MTUP2}: Our results suggest the two-step approach is a robust starting point for new languages
    \item \textbf{Decomposition is not always beneficial}: Avoid adding steps speculatively; validate on development data
    \item \textbf{Language-specific tuning matters}: Cross-linguistic transfer may require variant selection based on morphological complexity and resource availability
    \item \textbf{Computational efficiency}: LoRA-based fine-tuning on 7B models is practical and competitive for low-resource scenarios
\end{itemize}

\subsection{Future Work}

\begin{enumerate}
    \item \textbf{Multilingual scaling}: Extend evaluation to 5-10 typologically diverse languages to understand when decomposition helps
    \item \textbf{Larger models}: Evaluate with LLaMA 3, Mistral, and proprietary models to understand model-size dependencies
    \item \textbf{Domain transfer}: Test on news, biomedical, and social media texts to assess domain robustness
    \item \textbf{Hybrid approaches}: Combine MTUP with auxiliary objectives (dependency parsing, semantic role labeling) to improve performance
    \item \textbf{Error analysis}: Systematic analysis of failure modes and their relationship to decomposition choices
    \item \textbf{Few-shot in-context learning}: Incorporate gold examples into prompts to leverage in-context learning alongside fine-tuning
\end{enumerate}

\section{Conclusion}

We introduce Multi-Unified Prompting (MTUP), a framework for systematically exploring decomposition strategies in semantic parsing. This work makes two key contributions:

**First**, we present the **first neural AMR parser for Vietnamese**, establishing a baseline for future work on low-resource semantic parsing in morphologically complex languages.

**Second**, we provide empirical evidence that **decomposition effectiveness is language-dependent**, not universal. MTUP2 achieves 0.6494 SMATCH on Vietnamese by separating structure generation from variable binding, whereas direct end-to-end generation (B1, 0.6755) remains optimal for English. This cross-linguistic reversal is an important finding: one-size-fits-all decomposition strategies do not work.

**Key insight:** More decomposition steps do not always help. We find that error propagation and over-specification (too many intermediate constraints) can hurt performance. MTUP2 strikes an optimal balance for Vietnamese.

Rather than claiming a universal method, we contribute a systematic framework and empirical findings that guide decomposition design. Our results highlight the importance of language-aware prompt engineering for semantic parsing, especially in low-resource settings.

Future work should:
\begin{itemize}
    \item Extend to typologically diverse languages (5-10+) to understand when decomposition helps
    \item Investigate domain transfer (news, biomedical, formal text)
    \item Explore hybrid decomposition strategies
    \item Conduct fine-grained error analysis by semantic phenomenon
    \item Test larger models and in-context learning approaches
\end{itemize}

This work establishes Vietnamese AMR parsing as a new research direction and provides practical guidance for prompt engineering in multilingual semantic parsing.

\section*{Acknowledgments}

We thank [TODO: funding sources, annotation contributors, computational resources]. This work was partially supported by [TODO: grants, organizations].

\begin{thebibliography}{99}

\bibitem{banarescu2013abstract}
Banarescu, L., Bonial, C., Cai, S., Georgescu, M., Griffitt, K., Hermjakob, U., Knight, K., Koehn, P., Palmer, M., \& Schneider, N. (2013).
Abstract Meaning Representation for sembanking.
In \textit{Proceedings of the Linguistic Annotation Workshop}, pp. 178--186.

\bibitem{damonte2017incremental}
Damonte, M., Cohen, S. B., \& Satta, G. (2017).
An incremental parser for Abstract Meaning Representation.
In \textit{Proceedings of the 15th Conference of the European Chapter of the Association for Computational Linguistics}, pp. 536--546.

\bibitem{wang2015broad}
Wang, C., Xue, N., \& Pradhan, S. (2015).
A transition-based algorithm for AMR parsing.
In \textit{Proceedings of the 2015 Conference of the North American Chapter of the Association for Computational Linguistics}, pp. 366--375.

\bibitem{zhang2019amr}
Zhang, S., Ma, X., Duh, K., \& Van Durme, B. (2019).
AMR parsing as sequence-to-graph transduction.
In \textit{Proceedings of the 57th Annual Meeting of the Association for Computational Linguistics}, pp. 80--89.

\bibitem{naseem2019recurrent}
Naseem, T., Barzilay, R., \& Globerson, A. (2019).
Crossing the formality gap: Recurrent neural networks for informal social media parsing.
In \textit{Proceedings of the 2019 Conference of the North American Chapter of the Association for Computational Linguistics}, pp. 2219--2231.

\bibitem{jia2016semantic}
Jia, R., \& Liang, P. (2016).
Data recombination for neural semantic parsing.
In \textit{Proceedings of the 54th Annual Meeting of the Association for Computational Linguistics}, pp. 12--22.

\bibitem{guo2020towards}
Guo, D., Ren, S., Lu, S., Feng, Z., Tang, D., Liu, S., Zhou, L., Duan, N., Jiang, D., Zhou, M., \& others. (2021).
GraphCodeBERT: Pre-training code representations with data flow.
In \textit{International Conference on Learning Representations}.

\bibitem{bevilacqua2021spring}
Bevilacqua, M., Bisk, Y., \& Titov, I. (2021).
One SPRING to rule them both: Symmetric AMR semantic parsing and generation without complex graph transduction.
In \textit{Proceedings of the 59th Annual Meeting of the Association for Computational Linguistics}, pp. 6655--6668.

\bibitem{xie2023universal}
Xie, L., Gao, Y., Qin, B., \& Pan, S. J. (2023).
Instruction tuning for semantic parsing.
In \textit{Findings of the Association for Computational Linguistics: ACL 2023}.

\bibitem{li2023gpt}
Li, J., Constela, J., Liu, X., Bansal, M., \& Huang, Y. (2023).
GPT-3 and large language models for semantic parsing.
In \textit{Proceedings of the 2023 Conference on Empirical Methods in Natural Language Processing}.

\bibitem{wei2022chain}
Wei, J., Wang, X., Schuurmans, D., Bosma, M., Ichien, B., Xia, F., Chi, E., Le, Q. V., \& Zhou, D. (2022).
Emergent abilities of large language models.
arXiv preprint arXiv:2206.07682.

\bibitem{nye2021show}
Nye, M. I., Andreassen, A. J., Gur-Ari, G., Michalewski, H., Austin, J., Bieber, D., Dohan, D., Lewkowycz, A., Lewkowycz, E., Parisi, A., \& others. (2021).
Show your work: Improved reporting of experimental results.
In \textit{International Conference on Machine Learning}, pp. 8022--8032. PMLR.

\bibitem{shin2021constrained}
Shin, R., Lin, C. H., Thomson, S., Chen, B., Polozov, O., \& Choi, Y. (2021).
Constrained decoding for neural text generation from knowledge graphs.
In \textit{Findings of the Association for Computational Linguistics: EMNLP 2021}, pp. 1047--1057.

\bibitem{damonte2018cross}
Damonte, M., Cohen, S. B., \& Abend, O. (2018).
Cross-lingual abstract meaning representation parsing.
In \textit{Proceedings of the 2018 Conference of the North American Chapter of the Association for Computational Linguistics}, pp. 622--632.

\bibitem{xue2014not}
Xue, N., Li, J., \& Zhu, S. (2014).
Not an interlingua, but close: comparison of English-Chinese translations of eventuality types and event duration.
In \textit{Proceedings of the 9th International Conference on Language Resources and Evaluation (LREC'14)}, pp. 4076--4083.

\bibitem{nguyen2009vietnamese}
Nguyen, P. T., Vu, X. L., Tran, T. A., Phan, M. L., \& Ha, T. M. (2009).
Building a large syntactically-annotated corpus for Vietnamese.
In \textit{Proceedings of the Third Linguistic Annotation Workshop}, pp. 182--185.

\bibitem{vu2018vislp}
Vu, X. L., Miwa, M., Iwakura, T., Sakti, S., Asahara, M., Phan, M. L., Okazaki, N., \& Matsumoto, Y. (2018).
The VLSP shared task for Vietnamese NLP: Overview and results.
arXiv preprint arXiv:1807.11794.

\bibitem{chen2020low}
Chen, B., Shi, W., \& Ding, B. (2020).
Data augmentation for low-resource semantic parsing.
In \textit{Proceedings of the 2020 Conference on Empirical Methods in Natural Language Processing}, pp. 6045--6056.

\bibitem{artetxe2019cross}
Artetxe, M., Ruder, S., \& Yeniterzi, R. (2019).
Crosslingual language model pretraining.
In \textit{Advances in Neural Information Processing Systems}, pp. 7059--7069.

\bibitem{pfeiffer2020adapters}
Pfeiffer, J., Vulić, I., Gurevych, I., \& Schwenk, H. (2020).
MAD-X: An Adapter-Based Framework for Multi-Task Cross-Lingual Transfer.
In \textit{Proceedings of the 2020 Conference on Empirical Methods in Natural Language Processing}, pp. 7173--7192.

\bibitem{houlsby2019parameter}
Houlsby, N., Giurgiu, A., Jastrzebski, S., Morcos, A., de Freitas, N., \& Grangier, D. (2019).
Parameter-efficient transfer learning for NLP.
In \textit{International Conference on Machine Learning}, pp. 2790--2799. PMLR.

\bibitem{hu2021lora}
Hu, E. J., Shen, Y., Wallis, P., Allen-Zhu, Z., Li, Y., Wang, S., Wang, L., \& Chen, W. (2021).
LoRA: Low-Rank Adaptation of Large Language Models.
In \textit{International Conference on Learning Representations}.

\bibitem{zhang2023qlora}
Dettmers, T., Lewis, M., Belkada, Y., \& Zettlemoyer, L. (2023).
QLoRA: Efficient Finetuning of Quantized LLMs.
In \textit{Advances in Neural Information Processing Systems}.

\bibitem{wang2023lora}
Wang, S., Sun, Y., Xiang, Y., Wu, Z., Ding, S., Jing, N., \& Liu, M. (2023).
Scaling down to scale up: A guide to parameter-efficient fine-tuning.
arXiv preprint arXiv:2303.15647.

[TODO: Add more recent citations if needed]

\end{thebibliography}

\end{document}
